\section{Conclusiones}
\label{sec:conclusiones}

\begin{table}[htbp]
    \centering
    \caption{Conclusiones del an\'{a}lisis de flexibilidad de la l\'{i}nea SIM-014.}
    \label{tab:conclusiones}
    \setcounter{itemcount}{0}
    \begin{tabularx}{\textwidth}{@{}NX@{}}
        \toprule
        \multicolumn{1}{c}{\textbf{Item}} & \textbf{Conclusion} \\
        \midrule
        & La tuber\'{i}a de descarga del TMS (SIM-014) cumple los criterios de esfuerzo de ASME B31.3-2016 en los nueve casos de carga evaluados: la evaluaci\'{o}n de cumplimiento de c\'{o}digo del modelo arroj\'{o} CODE COMPLIANCE EVALUATION PASSED. \\
        & El esfuerzo de c\'{o}digo m\'{a}ximo es \num{64417.1}~\si{kPa} en el nodo 70, caso 6 (SUS) W+P2, con una relaci\'{o}n de utilizaci\'{o}n de 55.9~\% frente al admisible de \num{115142.4}~\si{kPa}; la condici\'{o}n que gobierna el dise\~{n}o es la sostenida (peso m\'{a}s presi\'{o}n), no el rango de expansi\'{o}n t\'{e}rmica. \\
        & El punto de mayor utilizaci\'{o}n corresponde a la tee no reforzada del nodo 70, por la concentraci\'{o}n de esfuerzo propia de este accesorio; aun as\'{i}, conserva un margen del 44.1~\% respecto al admisible. \\
        & Los rangos de desplazamiento t\'{e}rmico del sistema (m\'{a}ximo 39.2~\si{mm} en la componente $DZ$ del caso 8, EXP) se absorben con ratios de expansi\'{o}n inferiores a 0.2~\%, y las cargas en soportes son moderadas (carga vertical m\'{a}xima \num{1728}~\si{N} en el soporte del nodo 105 en operaci\'{o}n), por lo que la configuraci\'{o}n de soporter\'{i}a del isom\'{e}trico es adecuada. \\
        & La comparaci\'{o}n de las cargas en la conexi\'{o}n de descarga del TMS con los admisibles del fabricante del equipo queda pendiente como verificaci\'{o}n externa, por no disponerse de dichos admisibles a la fecha; el reporte de restricciones del job no incluye las cargas de los extremos del tramo. \\
        \bottomrule
    \end{tabularx}
\end{table}

En s\'{i}ntesis, la l\'{i}nea SIM-014 es t\'{e}cnicamente viable seg\'{u}n
ASME B31.3-2016 con una utilizaci\'{o}n m\'{a}xima de 55.9~\% del admisible.
No se identifican riesgos de sobre-esfuerzo, interferencias por
expansi\'{o}n ni sobrecarga de soportes, por lo que no se requieren acciones
correctivas sobre la configuraci\'{o}n analizada.
